
\documentclass{article}
\usepackage{colortbl}
\usepackage{makecell}
\usepackage{multirow}
\usepackage{supertabular}

\begin{document}

\newcounter{utterance}

\centering \large Interaction Transcript for game `cladder', experiment `full\_v1.5\_default', episode 8751 with qwen3.5{-}27b.
\vspace{24pt}

{ \footnotesize  \setcounter{utterance}{1}
\setlength{\tabcolsep}{0pt}
\begin{supertabular}{c@{$\;$}|p{.15\linewidth}@{}p{.15\linewidth}p{.15\linewidth}p{.15\linewidth}p{.15\linewidth}p{.15\linewidth}}
    \# & \multicolumn{2}{c}{Player} && \multicolumn{2}{c}{Game Master} \\
    \hline

    \theutterance \stepcounter{utterance}  
    & & & \multicolumn{4}{p{0.6\linewidth}}{
        \cellcolor[rgb]{0.9,0.9,0.9}{
            \makecell[{{p{\linewidth}}}]{
                \texttt{\tiny{[P1$\langle$GM]}}
                \texttt{You are an expert in causal inference. The following question is not a typical commonsense query, but rather a meticulously designed question created by a professor specializing in causal inference, intended to assess the students' mastery of the course content.} \\
\texttt{Imagine a self{-}contained, hypothetical world with only the following conditions, and without any unmentioned factors or causal relationships: Yupt has a direct effect on zupj. Zupj has a direct effect on muvq. For those who are not yupt, the probability of muvq is 25\%. For those who are yupt, the probability of muvq is 25\%. For those who are yupt, would it be more likely to see muvq if the individual was not yupt?} \\
\texttt{Start your answer with 'Yes' or 'No', followed by additional reasoning or evidence to support your explanation.} \\
            }
        }
    }
     \\ \\

    \theutterance \stepcounter{utterance}  
    & \multicolumn{4}{p{0.6\linewidth}}{
        \cellcolor[rgb]{0.561,0.737,0.561}{
            \makecell[{{p{\linewidth}}}]{
                \texttt{\tiny{[P1$\rangle$GM]}}
                \texttt{No} \\
\\ 
\texttt{The reasoning is as follows: The question asks whether an individual who is currently "yupt" would be more likely to see "muvq" if they were instead "not yupt." This is a counterfactual query comparing the probability of the outcome under two different states of the treatment variable (Yupt).} \\
\\ 
\texttt{According to the conditions provided:} \\
\texttt{1.  The probability of `muvq` given `yupt` is 25\% (\$P(muvq | yupt) = 0.25\$).} \\
\texttt{2.  The probability of `muvq` given `not yupt` is 25\% (\$P(muvq | \textbackslash{}neg yupt) = 0.25\$).} \\
\\ 
\texttt{Although the problem states there is a causal chain (\$Yupt \textbackslash{}to Zupj \textbackslash{}to Muvq\$), the observed marginal probabilities of the outcome `muvq` are identical regardless of the state of `yupt`. In a self{-}contained world with no unmentioned factors, if the probability of the outcome is the same in both the treated and untreated groups, the Average Treatment Effect (ATE) is zero.} \\
\\ 
\texttt{Therefore, changing the individual's state from `yupt` to `not yupt` would not increase the likelihood of `muvq`; the probability would remain exactly 25\%. Since 25\% is not greater than 25\%, it would not be "more likely" to see `muvq` under the counterfactual condition.} \\
            }
        }
    }
    & & \\ \\

    \theutterance \stepcounter{utterance}  
    & & & \multicolumn{2}{p{0.3\linewidth}}{
        \cellcolor[rgb]{0.9,0.9,0.9}{
            \makecell[{{p{\linewidth}}}]{
                \texttt{\tiny{[GM$|$GM]}}
                \texttt{No} \\
            }
        }
    }
    & & \\ \\

    \theutterance \stepcounter{utterance}  
    & & & \multicolumn{2}{p{0.3\linewidth}}{
        \cellcolor[rgb]{0.9,0.9,0.9}{
            \makecell[{{p{\linewidth}}}]{
                \texttt{\tiny{[GM$|$GM]}}
                \texttt{no} \\
            }
        }
    }
    & & \\ \\

    \theutterance \stepcounter{utterance}  
    & & & \multicolumn{2}{p{0.3\linewidth}}{
        \cellcolor[rgb]{0.9,0.9,0.9}{
            \makecell[{{p{\linewidth}}}]{
                \texttt{\tiny{[GM$|$GM]}}
                \texttt{game\_result = WIN} \\
            }
        }
    }
    & & \\ \\

\end{supertabular}
}

\end{document}
